\begin{description}
  \item[Corrections.] This section records substantive corrections, clarifications, and tracked revisions that affect the interpretation, notation, or results presented elsewhere in the book. When a later pass amends an earlier statement, the corrected form should be preferred, while the earlier form is retained here as part of the change history.

  \item[Version history.] The change log provides a compact record of the book's revision state, including the date or version identifier of each update, the scope of the modification, and a brief note on the affected material. This supports traceability across drafts and helps readers distinguish stable content from material that has been revised.

  \item[Limitations.] The book is intentionally selective in scope. Some topics are treated at the level needed for the main narrative rather than as exhaustive surveys, and some constructions are presented with assumptions that should be checked against the intended application domain. Any such limitations are noted here when they materially affect interpretation or reuse.

  \item[Known issues and open problems.] This log also tracks unresolved questions, gaps in coverage, and items deferred for future work. These may include proof obligations not yet discharged, implementation details awaiting validation, or extensions that remain open for further investigation.
\end{description}

The errata log is intended to remain a living record rather than a static appendix. As the manuscript evolves, entries should be updated to reflect corrected statements, newly identified limitations, and open problems that are relevant to readers who rely on the book as a technical reference. Where possible, each entry should identify the affected section or artifact and summarize the practical impact of the change.

For consistency with the rest of the book, future entries should be concise, versioned, and specific about scope. A useful entry format is: affected section or artifact, nature of the correction or limitation, version or date, and any downstream consequence for proofs, examples, or reproducibility artifacts. This keeps the log actionable for readers who need to determine whether a later revision changes a claim they depend on.
